\honest{Dissolving the Question}{Eliezer Yudkowsky}{2008}

Earlier I did not answer the riddle of the falling tree. I explained how the mind processes words, and I hope that at the end ``there was no question left---not even the feeling of a question.''

Many philosophers, especially amateur and ancient ones, share ``a dangerous instinct: If you give them a question, they try to answer it.'' Some weigh the arguments about free will and publish a verdict in a prestigious journal. Wiser ones recall that most philosophical disputes are about the meaning of a word, define the term precisely, and ask again. Wiser still, they argue that the idea is self-contradictory or meaningless for lack of testable consequences, and publish that. But proving that you are confused may not make you less confused, and proving a question meaningless may help no more than answering it. I name no philosopher. \nb{Wittgenstein wrote that the clarity he aimed at meant that philosophical problems ``should completely disappear.''}

But people's feeling of free will is a fact about human minds, and rejecting the idea does not explain the feeling: if free will does not exist, what goes on in the head of someone who thinks it does? ``This is not a rhetorical question!'' Likewise, noting that the tree's arguers anticipate nothing different and calling the argument pointless is correct, but does not say why they made the mistake. And mistakes reveal how a mind works. Our algorithms are how the world feels to us, and some things in the mind cut skew to the world.

My explanation is a network with a central unit, a useful shortcut that corresponds to nothing real, which stays undecided even after every observable feature is known. That unit feels like an unanswered question. ``(Metaphorically speaking. Human neurobiology is surely far more complex.)''

Once you understand ``in detail'' how your brain produces the feeling of the question, or better, the workings of Naive Bayes, ``then you're done,'' with no lingering confusion. Any vague dissatisfaction, or feeling of having been fast-talked, is a warning that you are not done. A thundering refutation of free will feels like a cheer for the home team, so you may not notice, or want to admit, that you cannot explain how each intuition arises.

An evolutionary story, such as that believers in free will outbred other hunter-gatherers, only argues that the brain produces an illusion, or explains why, not how; the illusion stays a brute fact. Likewise, guessing that people who called arguments meaningless lost status, so we evolved to argue about words, explains why; the network in ``Feel the Meaning'' explains how, breaking the confusion into pieces that are not confusing. I call my example of such a story ``a completely bogus explanation.'' \nb{The central unit is also a story, about mechanism instead of origin, and the post gives no more evidence for it than for the bogus one.}

Good hypotheses about mental algorithms are much harder than refuting a confusion, a different art altogether, so feel free to say: I can prove you wrong, but I cannot yet flowchart how your brain errs, so I am not done. Do not stop too early: it ``sometimes seems to me that at least 20\%'' of a skilled rationalist's effectiveness comes from that. I give no basis for the figure. The challenge is to notice even a little confusion, even with someone across from you insisting on free will and smirking.

When you are really done, ``you'll know you're done.'' It is ``an unmistakable feeling.'' Then: ``When you're done, you'll know you're done, but unfortunately the reverse implication does not hold.'' So the feeling can mislead you: ``Those who dream do not know they dream, but when you wake you know you are awake.'' My other test is being able to walk through the algorithm ``step by step.''

The homework: describe the mental algorithm that produces the debate about free will. Do not argue that people have free will or that they do not, and do not argue ``that free will is compatible with determinism, or not,'' or that the question is ill-posed or the concept contradictory or untestable, or tell an evolutionary story, or tie the concept to some bias; all of those explain why, not how. Write ``a stack trace'' of the mind's algorithms as they produce the intuitions behind the whole argument. It was one of the first real challenges I tried, and one of the easier ones. \nb{Three months later, in ``Thou Art Physics'' and ``Timeless Control,'' Yudkowsky gave an answer, and it argued that determinism is compatible with real control: one of the answers the homework had told readers not to give.}
